\subsection{Plurality with Elimination}


\content{
        \begin{example} Who wins the following election under the plurality
        method? \\
        \begin{tabular}{|l|c|c|c|c|c|c|}
        \hline
        & 3&4&4&6&2&1\\
        \hline
        1st & B&C&B&D&B&E\\
        \hline
        2nd & C&A&D&C&E&A\\
        \hline
        3rd & A&D&C&A&D&D\\
        \hline
        4th & D&B&A&E&C&B\\
        \hline
        5th & E&E&E&B&D&C\\
        \hline
        \end{tabular}
        \end{example}
}{% presentation
\vspace{4in}
}{% nopresentation
\vspace{2.5in}
}{% comments
Since the winner under plurality (B) does not have the majority, with the
plurality with eliminiation method we eliminate the last place candidate
(according to plurality) and recount the votes.
}

\content{
        \begin{example} Your turn! Use the Plurality with Elimination method to
        determine who wins. \\
        \begin{tabular}{|l|c|c|c|c|c|}
        \hline
        & 44 & 14 & 70 & 22 & 80 \\
        \hline
        1st & G&G&M&M&B\\
        \hline
        2nd & M&B&G&B&M\\
        \hline
        3rd & B&M&B&G&G\\
        \hline
        \end{tabular}
        \end{example}
}{% presentation
\vspace{3in}
}{% nopresentation
\vspace{2in}
}{% comments
$G$ is eliminated first, then $M$ has a majority. 
}

\content{
\begin{example} Use the plurality method with elimination to determine who wins
the election.\\
\begin{tabular}{|l|c|c|c|}
\hline
& 342 & 214 & 298 \\
\hline
1st & E&D&K\\
\hline
2nd & D&K&D\\
\hline
3rd & K&E&E\\
\hline
\end{tabular}

Is there a Condorcet Candidate? If so, did they win the election? 
\end{example}
}{% presentation
\vspace{3in}
}{% nopresentation
\vspace{2in}
}{% comments
$K$ wins under this method. But $D$ is a Condorcet Candidate. 
}

\content{
\begin{example} Who wins the following election under the plurality with
elimination method? \\
\begin{tabular}{|l|c|c|c|c|}
\hline
& 37 & 22 & 12 & 29 \\
\hline
1st & A&B&B&C\\
\hline
2nd & B&C&A&A\\
\hline
3rd & C&A&C&B\\
\hline
\end{tabular}
\end{example}
}{% presentation
\vspace{2in}
}{% nopresentation
\vspace{2in}
}{% comments
$C$ is eliminated first, then $A$ wins with the majority. 
}

\content{
Now, suppose that after officials have announced the results, the ballots were
accidentally destroyed before they could be certified, so the votes have to be
recast. Ten of the people whose preference was $B,A,C$ decide that they want to
now support the winner $A$, so they change their votes to $A,B,C$, all the other
votes remain the same. What are the results of the election now? \\
\begin{tabular}{|l|p{0.25in}|p{0.25in}|p{0.25in}|p{0.25in}|}
\hline
&  &  &  &  \\
\hline
1st & &&&\\
\hline
2nd & &&&\\
\hline
3rd & &&&\\
\hline
\end{tabular}
}{% presentation
\vspace{3in}
}{% nopresentation
\vspace{2in}
}{% comments
}


\content{
\begin{definition} (Monotonicity Criterion) If voters change their votes to
increase the preference of the winning candidate, it should not harm the
candidates chances of winning. 
\end{definition}
}{% presentation
}{% nopresentation
}{% comments
}

